% !TeX root = ../教材题库_全部条目_转写工作稿.tex
% Source: XJ-B1-C05-EX-P200-Y0177
% 第5章 / 5.5 三角函数模型的简单应用 / 教材 p.200 / PDF 第207页
% OCR未校对：公式与文字均可能有误；未完成数学TeX转写；禁止直接发给学生
% 原题及相邻原解截图（核对用）：../题目原图/XJ-B1-C05-EX-P200-Y0177-01.webp
% 本文件只在转写完成后作为题干使用；OCR 条目尚不是可靠题干。
\begin{bankitem}{5.5 例题 例题2}{第5章 / 5.5 三角函数模型的简单应用 / 教材 p.200 / PDF 第207页}{XJ-B1-C05-EX-P200-Y0177}{OCR未校对：公式与文字均可能有误；未完成数学TeX转写；禁止直接发给学生}
\begin{ocrdraft}
[教材p.200 / PDF第207页]
\par
00已知摩天轮的半径为60m，其中心距离地面70m，摩天轮做匀速转动，
\par
每30转一圈，摩天轮上点P的起始位置在最低点处．
\par
（1)试确定在时刻t(min)时，点P离地面的高度h；
\par
（2）在摩天轮转动的一圈内，点P离地面的高度超过100m的时间有多长？
\par
0
\par
/7
\par
一\}10m
\par
图5．5一2
\par
解（1）如图5．5一2，以摩天轮所在面为坐标平面，以摩天轮的中心点为原
\par
点0，轴和，轴分别平行和垂直于地平面，建立平面直角坐标系．
\par
点P的初始位置在最低点，设点P从最低点沿逆时针方向匀速转动，OP在时
\par
间t(min)内所转过的角度为“．'一玉／，因此OP与的夹角为一玉十匹／，于是，
\par
30
\par
215
\par
点P的纵坐标3，一60sin
\par
—60cos一一t.
\par
因此，点P离地面的高度h一70十3，一70一60cos一／（m)．
\par
（2）根据题意，由于0K30，则严[0，2丌]
\par
由
\par
70一60cos．／>100，
\par
15
\par
解得
\par
3
\par
15
\par
从而
\par
10<十<20\ensuremath{\cdot}
\par
因此，在摩天轮转动的一圈内，点P离地面的高度超过1佣m的时间有10min.
\par
三角函数作为描述现实世界中周期现象的重要数学模型有着广泛的应用．正如
\par
章前语的描述：潮涨潮落、月圆月缺、四季交替等是自然界中按一定的规律周而复
\par
始出现的现象“““这都说明周期变化的现象在现实生活中比比皆是．
\end{ocrdraft}
\end{bankitem}
